\documentclass[11pt,a4paper]{article}
\usepackage{fontspec,amsmath,amssymb,amsthm,microtype}
\setmainfont{lmroman10-regular.otf}[BoldFont=lmroman10-bold.otf,ItalicFont=lmroman10-italic.otf,BoldItalicFont=lmroman10-bolditalic.otf]
\setmonofont{lmmono10-regular.otf}[Scale=MatchLowercase]
\usepackage[margin=25mm]{geometry}
\usepackage[unicode,hidelinks]{hyperref}
\setlength{\emergencystretch}{3em}
\setlength{\parskip}{0.4em}
\setlength{\parindent}{0pt}
\allowdisplaybreaks[2]
\raggedbottom
\title{\texorpdfstring{Construct the L\ensuremath{{}^{2}} Langevin equation from Brownian-\allowbreak{}coordinate laws}{Construct the L² Langevin equation from Brownian-coordinate laws}}
\author{LangevinLean mathematical exposition}
\date{}
\begin{document}\maketitle
\noindent Source module: \texttt{Unadjusted\allowbreak{}Langevin/\allowbreak{}Brownian\allowbreak{}Langevin.\allowbreak{}lean}. The accompanying HTML pairs this mathematical exposition with the exact Lean statements and records the scope of the formalization.


Coordinate pre-Brownian laws produce the continuous \(L^2\) additive noise needed by the deterministic integral-equation existence theorem. The exact noise norm is \(\sqrt{2d}\sqrt t\).


\section{\texorpdfstring{brownian\allowbreak{}Langevin\allowbreak{}Noise}{brownianLangevinNoise}}


Define \(N_a(t)=\sqrt2([B_{a+t}]-[B_a])\) in \(L^2\), using the real-time Brownian extension. The subtraction fixes the input at zero at local time zero, and the \(\sqrt2\) factor gives the overdamped Langevin diffusion normalization.



{\raggedright\noindent Lean declarations: \texttt{brownian\allowbreak{}Langevin\allowbreak{}Noise}.\par}

\section{\texorpdfstring{brownian\allowbreak{}Langevin\allowbreak{}Noise continuous}{brownianLangevinNoise continuous}}


The real-time Brownian map is continuous in \(L^2\) from its Gaussian increment moments. Composition with \(t\mapsto a+t\), subtraction of the fixed class \([B_a]\), and scalar multiplication preserve continuity. Thus the additive input satisfies the deterministic Picard theorem's continuity requirement.



{\raggedright\noindent Lean declarations: \texttt{brownian\allowbreak{}Langevin\allowbreak{}Noise\_\allowbreak{}continuous}.\par}

\section{\texorpdfstring{brownian\allowbreak{}Langevin\allowbreak{}Noise zero}{brownianLangevinNoise zero}}


At local time zero both Brownian terms are \([B_a]\), so they cancel and their scalar multiple is zero. Lean unfolds the definition and uses additive identities; no exceptional-set statement is needed because this is equality in \(L^2\).



{\raggedright\noindent Lean declarations: \texttt{brownian\allowbreak{}Langevin\allowbreak{}Noise\_\allowbreak{}zero}.\par}

\section{\texorpdfstring{brownian\allowbreak{}Langevin\allowbreak{}Noise norm}{brownianLangevinNoise norm}}


For \(a,t\ge0\), apply the exact real-time Brownian increment moment formula. It gives \(\|N_a(t)\|_2=\sqrt{2d}\sqrt t\), with \(d=|\iota|\). The Gaussian laws prove this value; it is not an additional noise-bound hypothesis.



{\raggedright\noindent Lean declarations: \texttt{brownian\allowbreak{}Langevin\allowbreak{}Noise\_\allowbreak{}norm}.\par}

\section{\texorpdfstring{exists lp langevin integral of brownian small}{exists lp langevin integral of brownian small}}


Supply the continuous, zero-starting \(N_a\) to \texttt{exists\_\allowbreak{}lp\_\allowbreak{}langevin\_\allowbreak{}integral\_\allowbreak{}small}. It constructs \(Y_t-Y_0=-\int_0^t\operatorname{lpDrift}(\nabla U)(Y_s)ds+N_a(t)\) for \(Lh\le1/2\). Attach the exact noise norm theorem to the existential package. Neither the integral equation nor its noise estimate is assumed.



{\raggedright\noindent Lean declarations: \texttt{exists\_\allowbreak{}lp\_\allowbreak{}langevin\_\allowbreak{}integral\_\allowbreak{}of\_\allowbreak{}brownian\_\allowbreak{}small}.\par}

\section{\texorpdfstring{exists lp langevin integral of brownian}{exists lp langevin integral of brownian}}


Use the unrestricted finite-horizon \(L^2\) existence theorem with the same constructed Brownian input. The resulting solution is continuous in \(L^2\), starts at the prescribed random variable, and satisfies the actual gradient integral equation. The exact \(\sqrt{2d}\sqrt t\) noise formula is included as a proved property.



{\raggedright\noindent Lean declarations: \texttt{exists\_\allowbreak{}lp\_\allowbreak{}langevin\_\allowbreak{}integral\_\allowbreak{}of\_\allowbreak{}brownian}.\par}

\end{document}
